\chapter{Campagnes de mesure d'autonomie}
\label{ann:autonomie}

Cette annexe rassemble les relevés bruts et les grandeurs dérivées des trois campagnes d'autonomie menées sur le prototype. Les valeurs synthétisées ici sont reprises et commentées dans le chapitre de résultats.

\section{Protocole}
\label{ann:autonomie:protocole}

Chaque campagne consiste en une décharge complète depuis une batterie pleine jusqu'à l'arrêt de l'appareil, sans intervention de l'utilisateur et sans alimentation externe connectée. Le service de mesure d'autonomie embarqué dans le firmware échantillonne périodiquement la jauge \gls{soc} MAX17260 et écrit une ligne de relevé sur la carte microSD. Les trois campagnes ont été exécutées avec la même version de firmware, identifiée par le condensé de commit \texttt{878e43d}, sur le même exemplaire de carte et la même cellule, afin que les écarts observés soient attribuables au mode de fonctionnement et non à une variation de plateforme.

Les trois modes évalués sont la veille, la lecture filaire sur le jack 3,5 mm et la lecture Bluetooth A2DP. En veille, l'appareil est allumé, l'écran est actif et aucun flux audio n'est produit. En lecture filaire et en lecture Bluetooth, le même fichier WAV 44,1 kHz est lu en boucle depuis la carte microSD. La fréquence du processeur est fixée à 160 MHz et l'interface microSD est cadencée à 40 MHz dans les trois cas.

Le courant et la tension proviennent de la jauge MAX17260 et non d'un instrument externe. Cette mesure interne est celle dont dispose le produit fini, ce qui la rend représentative de l'information réellement exploitable par le firmware, mais elle hérite de l'incertitude propre au circuit de jauge et de son shunt. Les valeurs absolues doivent donc être lues avec la réserve correspondante, les comparaisons relatives entre modes restant quant à elles pleinement valides puisque toutes issues du même dispositif de mesure.

\section{Format des relevés}
\label{ann:autonomie:format}

Chaque fichier de campagne est un CSV précédé d'un en-tête, qui consigne les conditions d'exécution. Le tableau \ref{tab:auton_colonnes} décrit les colonnes de données. La convention de signe est celle de la jauge, un courant négatif correspondant à une décharge. Toutes les grandeurs présentées dans la suite de cette annexe utilisent la convention inverse, un courant positif désignant le courant consommé par l'appareil.

\begin{table}[H]
    \centering
    \caption{Colonnes des fichiers de relevé d'autonomie}
    \label{tab:auton_colonnes}
    \begin{tabular}{lll}
        \toprule
        Colonne                & Unité & Description                                 \\
        \midrule
        \texttt{t\_s}          & s     & Temps écoulé depuis le début de la campagne \\
        \texttt{voltage\_mV}   & mV    & Tension de la cellule mesurée par la jauge  \\
        \texttt{soc\_pct}      & \%    & État de charge estimé par la jauge          \\
        \texttt{current\_mA}   & mA    & Courant de cellule, négatif en décharge     \\
        \texttt{bt}            & -     & Indicateur d'activité de la pile Bluetooth  \\
        \texttt{heap\_min\_kb} & kio   & Plancher de tas observé depuis le démarrage \\
        \bottomrule
    \end{tabular}
\end{table}

\section{Synthèse des trois campagnes}
\label{ann:autonomie:synthese}

Le tableau \ref{tab:auton_synthese} rassemble les grandeurs dérivées des trois relevés. Le courant moyen et la puissance moyenne sont obtenus par intégration trapézoïdale sur l'axe temporel, puis division par la durée totale, ce qui les rend insensibles à la période d'échantillonnage retenue pour chaque campagne. La charge et l'énergie consommées résultent des mêmes intégrales, exprimées respectivement en mAh et en mWh.

\begin{table}[H]
    \centering
    \caption{Synthèse des trois campagnes d'autonomie}
    \label{tab:auton_synthese}
    \begin{tabular}{lrrrrrr}
        \toprule
        Mode              & Durée & $I_\text{moy}$ & $I_\text{min}$ / $I_\text{max}$ & $P_\text{moy}$ & $Q$   & $E$   \\
                          & [h]   & [mA]           & [mA]                            & [mW]           & [mAh] & [mWh] \\
        \midrule
        Veille            & 13,46 & 62             & 36 / 96                         & 236            & 834   & 3175  \\
        Lecture jack      & 9,24  & 96             & 76 / 137                        & 366            & 890   & 3381  \\
        Lecture Bluetooth & 5,81  & 156            & 120 / 211                       & 586            & 903   & 3405  \\
        \bottomrule
    \end{tabular}
\end{table}

Les trois campagnes se terminent sur une tension de cellule comprise entre 3289 mV et 3306 mV, alors que le seuil de coupure du superviseur XC6120 est situé à 3,2 V. L'arrêt n'est donc pas provoqué par le superviseur mais par l'atteinte du plancher d'état de charge signalé par la jauge, ce qui confirme que la protection contre la décharge profonde reste une sécurité de dernier recours et non le mécanisme d'arrêt nominal.

La charge restituée s'échelonne de 834 mAh à 903 mAh, à comparer à la capacité nominale de 950 mAh annoncée par le fabricant de la cellule \cite{jauch_LP523450JU}. Cette valeur nominale est toutefois définie pour une décharge menée jusqu'à 3,0 V, alors que les campagnes s'interrompent aux alentours de 3,3 V. Une fraction de l'écart correspond donc simplement à la portion de décharge non exploitée sous ce seuil, et non à un déficit de la cellule.

Le taux de décharge ne peut en revanche pas expliquer la dispersion résiduelle entre les trois campagnes. La capacité nominale de décharge de la cellule est spécifiée à 0,2 C, soit 190 mA, alors que les trois modes consomment respectivement 62 mA, 96 mA et 156 mA, c'est-à-dire moins que la condition de référence. Un régime plus doux tend à augmenter la capacité restituée, or c'est la campagne de veille, la moins exigeante des trois, qui restitue le moins de charge. L'écart de 8 \% entre les campagnes est donc à attribuer à l'incertitude d'estimation de la jauge et à l'échantillonnage instantané du courant plutôt qu'à un effet de charge, ce qui fixe l'ordre de grandeur de la confiance à accorder aux valeurs absolues de charge.

Le plancher de tas relevé varie de 3871 kio à 3921 kio selon la campagne et reste strictement constant à l'intérieur de chaque campagne après les premières minutes de fonctionnement. Sur une décharge complète de 13 h en veille, cette stabilité constitue un indice de l'absence de fuite mémoire dans le régime permanent, l'écart entre campagnes correspondant simplement à l'empreinte des sous-systèmes activés dans chaque mode.

\section{Comparaison des modes}
\label{ann:autonomie:comparaison}

Le tableau \ref{tab:auton_comparaison} exprime chaque mode de lecture relativement au mode veille, qui constitue le plancher de consommation de la plateforme. Le surcoût en courant isole ainsi la consommation propre à la fonction évaluée, indépendamment du coût fixe du processeur, de l'écran et des périphériques toujours actifs.

\begin{table}[H]
    \centering
    \caption{Surcoût des modes de lecture par rapport à la veille}
    \label{tab:auton_comparaison}
    \begin{tabular}{lrrrr}
        \toprule
        Mode              & $\Delta I$ & $\Delta P$ & $I / I_\text{veille}$ & Perte d'autonomie \\
                          & [mA]       & [mW]       &                       & [\%]              \\
        \midrule
        Lecture jack      & +34        & +130       & 1,55                  & 31,3              \\
        Lecture Bluetooth & +94        & +351       & 2,51                  & 56,9              \\
        \bottomrule
    \end{tabular}
\end{table}

La lecture filaire ajoute 34 mA au plancher de veille. Ce surcoût regroupe la lecture de la carte microSD, le décodage et la mise en forme des échantillons, l'horloge \gls{i2s}, ainsi que la consommation du \gls{DAC} PCM5242 et de l'amplificateur MAX97220 débitant dans la charge du casque. Il porte l'autonomie de 13,5 h à 9,2 h, soit une perte de 31 \%.

La lecture Bluetooth ajoute 94 mA, soit un surcoût presque trois fois supérieur à celui de la chaîne filaire. La différence de 59 mA entre les deux modes de lecture correspond à l'écart entre le coût de l'émetteur radio et de l'encodage SBC d'une part, et le coût de la chaîne analogique désactivée dans ce mode d'autre part. L'autonomie tombe à 5,8 h, soit une perte de 57 \% par rapport à la veille. Le Bluetooth constitue ainsi le poste de consommation dominant de l'appareil.

Le réglage de volume retenu diffère entre les deux campagnes de lecture, pour des raisons propres à chaque mode. En lecture filaire, le volume est fixé à 50~\%, afin de garantir la reproductibilité de la campagne. Ce point de fonctionnement détermine directement la puissance délivrée au casque et donc le courant consommé par l'étage analogique, ce qui en fait un paramètre déterminant du résultat.

En A2DP, le volume est fixé à sa valeur minimale non nulle. Ce choix ne correspond pas à une condition d'écoute mais à une contrainte expérimentale, la plupart des récepteurs Bluetooth basculant automatiquement en veille après une période prolongée sans signal utile, ce qui aurait interrompu la campagne avant la décharge complète. Un flux de faible amplitude mais continu maintient le récepteur actif pendant toute la durée du test. Ce réglage n'introduit aucun biais sur la mesure, puisque l'encodage \gls{sbc} opère à débit fixe et que la consommation de l'émetteur ne dépend pas de l'amplitude des échantillons transmis.

\section{Relevés temporels}
\label{ann:autonomie:releves}

Les figures suivantes présentent, pour chaque campagne, l'évolution de l'état de charge, du courant consommé et de la tension de cellule en fonction du temps. Le courant est tracé à la fois brut et lissé par moyenne glissante, la trace brute rendant visible la dispersion d'échantillonnage et la trace lissée la tendance de consommation.

\fig[H, width=1\textwidth]{Campagne de veille, décharge complète en 13,46 h}{auton_idle.png}

\fig[H, width=1\textwidth]{Campagne de lecture filaire sur le jack 3,5 mm, décharge complète en 9,24 h}{auton_jack.png}

\fig[H, width=1\textwidth]{Campagne de lecture Bluetooth A2DP, décharge complète en 5,81 h}{auton_bt.png}

La figure \ref{auton_comparaison.png} superpose les trois campagnes sur un axe temporel commun. La linéarité des courbes d'état de charge sur la quasi-totalité de la décharge traduit un régime de consommation stable dans les trois modes, sans dérive progressive qui trahirait une accumulation de ressources ou une dégradation du fonctionnement au fil des heures.

\fig[H, width=1\textwidth]{Comparaison des trois campagnes, état de charge et courant consommé}{auton_comparaison.png}

La décroissance de la tension suit dans les trois cas la courbe caractéristique d'une cellule lithium-ion, avec un plateau étendu entre 4,0 V et 3,7 V puis une chute rapide en fin de décharge. Le décalage vertical entre les trois courbes de tension au même état de charge correspond à la chute ohmique sur la résistance interne de la cellule, d'autant plus marquée que le courant est élevé, ce qui explique que la campagne Bluetooth se termine sur la tension la plus basse des trois.